\paragraph{算术自然密度测度、可数 prime-register solenoid 与有限维可见性的统一塌缩}
在黄金地址柱测度的对数级质量律、无穷 prime-register solenoid 的有限可见因子化，以及联合有限证书不完备性已经建立之后，还可进一步抽取出下列五条结构后果。

\begin{theorem}[算术自然密度诱导的黄金柱测度不是任何连续势的 Gibbs 态]\label{thm:xi-time-part60f2-zg-density-measure-non-gibbs}
\leanverified{paper\_xi\_time\_part60f2\_zg\_density\_measure\_non\_gibbs}
沿用定理 \ref{thm:xi-time-part60f-zg-cylinder-measure-local-dimension-zero} 的记号，记
\[
\sigma:\Sigma_{\mathrm{ZG}}\to \Sigma_{\mathrm{ZG}},
\qquad
\sigma(z_1,z_2,z_3,\dots)=(z_2,z_3,\dots).
\]
则自然密度诱导的柱测度
\[
\nu_{\mathrm{ZG}}
\]
不是 \((\Sigma_{\mathrm{ZG}},\sigma)\) 上任何连续势
\[
\phi\in C(\Sigma_{\mathrm{ZG}})
\]
的 Gibbs 测度。
\end{theorem}
\begin{proof}
对每个 \(n\ge 1\)，记
\[
[0^n]:=[(0,\dots,0)]\subset \Sigma_{\mathrm{ZG}}.
\]
由定理 \ref{thm:xi-time-part60f-zg-zero-prefix-log-density} 与定理 \ref{thm:xi-time-part60f-zg-cylinder-measure-local-dimension-zero}，
\[
\nu_{\mathrm{ZG}}([0^n])
=
\frac{\mathrm{dens}(\mathrm{ZG}[0^n])}{\delta_{\mathrm{ZG}}}
\sim
\frac{e^{-\gamma}}{\delta_{\mathrm{ZG}}\log p_n}.
\]
特别地，
\[
\nu_{\mathrm{ZG}}([0^n])\to 0,
\]
并且该衰减仅为次指数级。

反设 \(\nu_{\mathrm{ZG}}\) 是某个连续势 \(\phi\in C(\Sigma_{\mathrm{ZG}})\) 的 Gibbs 测度。则按 Gibbs 定义，存在常数 \(C\ge 1\) 使对一切 \(n\ge 1\) 都有
\[
C^{-1}
\le
\nu_{\mathrm{ZG}}([0^n])
\exp\!\bigl(nP(\phi)-S_n\phi(0^\infty)\bigr)
\le
C,
\]
其中 \(0^\infty=(0,0,\dots)\in\Sigma_{\mathrm{ZG}}\) 为全零固定点。
由于
\[
S_n\phi(0^\infty)=n\phi(0^\infty),
\]
便得到
\[
\nu_{\mathrm{ZG}}([0^n])\asymp \exp\!\bigl(-n(P(\phi)-\phi(0^\infty))\bigr).
\]
若 \(P(\phi)<\phi(0^\infty)\)，则右端随 \(n\) 指数增长，与 \(\nu_{\mathrm{ZG}}([0^n])\le 1\) 矛盾。
若 \(P(\phi)=\phi(0^\infty)\)，则 \(\nu_{\mathrm{ZG}}([0^n])\) 有正下界，这与 \(\nu_{\mathrm{ZG}}([0^n])\to 0\) 矛盾。
若 \(P(\phi)>\phi(0^\infty)\)，则 \(\nu_{\mathrm{ZG}}([0^n])\) 必指数衰减；而
\[
\nu_{\mathrm{ZG}}([0^n])\sim \frac{e^{-\gamma}}{\delta_{\mathrm{ZG}}\log p_n}
\]
仅按 \((\log p_n)^{-1}\) 衰减。由素数定理 \(p_n\sim n\log n\) 可知
\[
\log\log p_n=o(n),
\]
故这不可能与任意指数律等价。三种情形皆矛盾，因此 \(\nu_{\mathrm{ZG}}\) 不可能是任何连续势的 Gibbs 测度。证毕。
\end{proof}

\begin{theorem}[可数 prime-register solenoid 之间连续态射的完全分类]\label{thm:xi-time-part60f2-countable-solenoid-morphism-classification}
\leanverified{paper\_xi\_time\_part60f2\_countable\_solenoid\_morphism\_classification}
设 \(S,T\subset\mathbb P\) 为可数素数集，并记
\[
G_S:=\ZZ[S^{-1}],
\qquad
G_T:=\ZZ[T^{-1}],
\qquad
\Sigma_S:=\widehat{G_S},
\qquad
\Sigma_T:=\widehat{G_T}.
\]
则存在自然同构
\[
\operatorname{Hom}_{\ZZ}(G_S,G_T)\cong
\begin{cases}
0,& S\not\subseteq T,\\[1mm]
G_T,& S\subseteq T,
\end{cases}
\]
从而对偶化得到
\[
\operatorname{Hom}_{\mathrm{cts}}(\Sigma_T,\Sigma_S)\cong
\begin{cases}
0,& S\not\subseteq T,\\[1mm]
G_T,& S\subseteq T.
\end{cases}
\]
\end{theorem}
\begin{proof}
先设 \(S\subseteq T\)。
任取 \(x\in G_T\)，定义
\[
m_x:G_S\to G_T,
\qquad
y\mapsto xy.
\]
由于 \(G_S,G_T\subset\QQ\)，且 \(y\in G_S\) 的分母素数仅来自 \(S\subseteq T\)，故 \(xy\in G_T\)，从而 \(m_x\) 为良定义的加法同态。

反过来，任取 \(\phi\in \operatorname{Hom}_{\ZZ}(G_S,G_T)\)，并记
\[
x:=\phi(1)\in G_T.
\]
对任意 \(y=a/b\in G_S\)（其中 \(a\in\ZZ\)，\(b\) 的全部素因子属于 \(S\)），有
\[
b\,\phi(y)=\phi(by)=\phi(a)=a\,\phi(1)=ax.
\]
由于 \(G_T\subset\QQ\) 无挠，故
\[
\phi(y)=\frac{a}{b}x=xy=m_x(y).
\]
因此 \(\phi=m_x\)，从而
\[
\operatorname{Hom}_{\ZZ}(G_S,G_T)\cong G_T.
\]

再设 \(S\not\subseteq T\)，取 \(p\in S\setminus T\)。
任取 \(\phi\in\operatorname{Hom}_{\ZZ}(G_S,G_T)\)，并记 \(x:=\phi(1)\)。
由于 \(p^{-n}\in G_S\) 对每个 \(n\ge 1\) 都成立，故
\[
x=\phi(1)=p^n\phi(p^{-n})\in p^nG_T
\qquad(n\ge 1).
\]
于是
\[
x\in \bigcap_{n\ge 1}p^nG_T.
\]
若 \(u\in G_T\setminus\{0\}\)，则可写成
\[
u=\frac{a}{\prod_{q\in T}q^{m_q}}
\]
其中 \(a\in\ZZ\setminus\{0\}\)，且分母与 \(p\) 互素。
若 \(u\in p^nG_T\)，则把两边写成既约分式可知 \(p^n\mid a\)。
这不可能对所有 \(n\) 同时成立，故
\[
\bigcap_{n\ge 1}p^nG_T=\{0\}.
\]
于是 \(x=0\)。
再对任意 \(y=a/b\in G_S\) 有
\[
b\,\phi(y)=\phi(a)=a\,\phi(1)=0,
\]
而 \(G_T\) 无挠，故 \(\phi(y)=0\)。
因此
\[
\operatorname{Hom}_{\ZZ}(G_S,G_T)=0.
\]

最后，Pontryagin 对偶给出反变等价
\[
\operatorname{Hom}_{\mathrm{cts}}(\Sigma_T,\Sigma_S)\cong \operatorname{Hom}_{\ZZ}(G_S,G_T),
\]
故连续态射的分类随即成立。证毕。
\end{proof}

\begin{theorem}[无穷 prime-register solenoid 的有限维连续表示只看见有限个素数轴]\label{thm:xi-time-part60f2-finite-dimensional-representation-finite-prime-factorization}
\leanverified{paper\_xi\_time\_part60f2\_finite\_dimensional\_representation\_finite\_prime\_factorization}
设 \(S\subset\mathbb P\) 为可数无穷素数集，并记
\[
\Sigma_S:=\widehat{\ZZ[S^{-1}]}.
\]
任取连续酉表示
\[
\rho:\Sigma_S\to U(d).
\]
则存在有限子集
\[
S_0\Subset S
\]
与连续酉表示
\[
\rho_0:\Sigma_{S_0}\to U(d)
\]
使得
\[
\rho=\rho_0\circ \pi_{S\to S_0}.
\]
因此余有限尾核
\[
K_{S\setminus S_0}:=\prod_{p\in S\setminus S_0}\ZZ_p
\]
满足
\[
K_{S\setminus S_0}\subseteq \ker\rho.
\]
此外，\(\rho(\Sigma_S)\) 的 Lie 秩至多为 \(1\)；更强地，
\[
\rho(\Sigma_S)
\]
只能是平凡群或一维圆群 \(\TT\)。
\end{theorem}
\begin{proof}
由于 \(\Sigma_S\) 为紧 Abel 群，\(\rho(\Sigma_S)\) 是一族彼此交换的酉矩阵，故可同时酉对角化。
于是存在 \(U\in U(d)\) 与连续特征
\[
\chi_1,\dots,\chi_d:\Sigma_S\to \TT
\]
使得对每个 \(\xi\in\Sigma_S\) 都有
\[
U\rho(\xi)U^{-1}
=
\operatorname{diag}\bigl(\chi_1(\xi),\dots,\chi_d(\xi)\bigr).
\]
由 Pontryagin 对偶，
\[
\operatorname{Hom}_{\mathrm{cts}}(\Sigma_S,\TT)\cong G_S:=\ZZ[S^{-1}],
\]
故每个 \(\chi_j\) 对应于某个 \(x_j\in G_S\)。
每个 \(x_j\) 的分母只涉及有限多个素数。
令 \(S_0\subset S\) 为所有 \(x_j\) 的分母中出现之素数的并集，则 \(S_0\) 有限，且 \(x_j\in G_{S_0}\) 对所有 \(j\) 都成立。
于是每个 \(\chi_j\) 都经由
\[
\pi_{S\to S_0}:\Sigma_S\twoheadrightarrow \Sigma_{S_0}
\]
因子化，从而整个 \(\rho\) 也经由 \(\pi_{S\to S_0}\) 因子化。

由推论 \ref{cor:xi-time-part60e-infinite-solenoid-finite-visible-factorization} 的短正合列，
\[
0\to \prod_{p\in S\setminus S_0}\ZZ_p\to \Sigma_S\xrightarrow{\ \pi_{S\to S_0}\ }\Sigma_{S_0}\to 0,
\]
立得
\[
K_{S\setminus S_0}\subseteq \ker\rho.
\]

最后分析像的 Lie 秩。
在对角化后，\(\rho\) 可视为映射
\[
\Sigma_S\to \TT^d.
\]
其对偶离散映射为
\[
\rho^\vee:\ZZ^d\to G_S,
\qquad
e_j\mapsto x_j.
\]
因此
\[
\widehat{\rho(\Sigma_S)}\cong \operatorname{im}(\rho^\vee)=:H
=
\langle x_1,\dots,x_d\rangle\subset G_S\subset\QQ.
\]
而 \(\QQ\) 的任意有限生成子群都循环：取 \(x_1,\dots,x_d\) 的一个公共分母 \(m\)，则
\[
H\subseteq \frac1m\ZZ,
\]
故 \(H\) 是 \((1/m)\ZZ\) 的有限生成子群，从而必为零群或无限循环群。
因此其 Pontryagin 对偶
\[
\rho(\Sigma_S)\cong \widehat H
\]
只能是平凡群或 \(\TT\)。
特别地，\(\rho(\Sigma_S)\) 的 Lie 秩至多为 \(1\)。证毕。
\end{proof}

\begin{corollary}[无穷 prime-register 终对象不存在忠实有限维 Lie 模型]\label{cor:xi-time-part60f2-no-faithful-finite-dimensional-lie-model}
\leanverified{paper\_xi\_time\_part60f2\_no\_faithful\_finite\_dimensional\_lie\_model}
在定理 \ref{thm:xi-time-part60f2-finite-dimensional-representation-finite-prime-factorization} 的设定下，\(\Sigma_S\) 不存在忠实的有限维连续酉表示。
更一般地，不存在到任何有限维紧 Lie 群的忠实连续商实现。
\end{corollary}
\begin{proof}
若 \(\rho:\Sigma_S\to U(d)\) 忠实，则由定理 \ref{thm:xi-time-part60f2-finite-dimensional-representation-finite-prime-factorization}，存在有限 \(S_0\Subset S\) 使
\[
K_{S\setminus S_0}=\prod_{p\in S\setminus S_0}\ZZ_p
\subseteq \ker\rho.
\]
由于 \(S\setminus S_0\) 非空，\(K_{S\setminus S_0}\) 为非平凡闭子群，从而 \(\ker\rho\neq \{0\}\)，矛盾。

若进一步存在一个到有限维紧 Lie 群的忠实连续商实现
\[
q:\Sigma_S\twoheadrightarrow G,
\qquad
\ker q=\{0\},
\]
则由 Peter--Weyl 定理，\(G\) 存在某个忠实有限维酉表示
\[
\iota:G\hookrightarrow U(d).
\]
于是
\[
\iota\circ q:\Sigma_S\to U(d)
\]
给出忠实有限维连续酉表示，与上一步矛盾。故所述 Lie 模型亦不存在。证毕。
\end{proof}

\begin{theorem}[有限半径 Jensen 证书与有限维连续表示的统一塌缩]\label{thm:xi-time-part60f2-finite-radius-finite-dimensional-unified-collapse}
设 \(S\subset\mathbb P\) 为可数无穷素数集，固定
\[
0<\rho_\star<1,
\]
并取任意有限维连续酉表示
\[
\varrho:\Sigma_S\to U(d).
\]
则存在一列对象对
\[
\{(F_n,\xi_n)\}_{n\ge 1},
\]
其中每个 \(F_n\) 都是单位圆盘上的有限 Blaschke 乘积，且 \(\xi_n\in \Sigma_S\)，满足
\begin{enumerate}
  \item \(\mathcal J_{\rho_\star}(F_n)\) 全部相同；
  \item \(\varrho(\xi_n)\) 全部相同；
  \item \(F_n\) 两两不同；
  \item \(\xi_n\) 两两不同。
\end{enumerate}
因此，任何把有限半径 Jensen 剖面与有限维连续算术观测并合的有限可见证书都不可能完备。
\end{theorem}
\begin{proof}
由定理 \ref{thm:xi-time-part60f2-finite-dimensional-representation-finite-prime-factorization}，存在有限子集 \(S_0\Subset S\) 与连续酉表示
\[
\varrho_0:\Sigma_{S_0}\to U(d)
\]
使
\[
\varrho=\varrho_0\circ \pi_{S\to S_0}.
\]
再由定理 \ref{thm:xi-time-part60e-finite-radius-finite-prime-certificate-incompleteness}，存在一列对象对
\[
\{(F_n,\xi_n)\}_{n\ge 1}
\]
使得
\[
\mathcal J_{\rho_\star}(F_n)
=
\mathcal J_{\rho_\star}(F_1),
\qquad
\pi_{S\to S_0}(\xi_n)
=
\pi_{S\to S_0}(\xi_1),
\]
并且 \(F_n\) 两两不同、\(\xi_n\) 两两不同。
于是对所有 \(n\ge 1\)，
\[
\varrho(\xi_n)
=
\varrho_0\!\bigl(\pi_{S\to S_0}(\xi_n)\bigr)
=
\varrho_0\!\bigl(\pi_{S\to S_0}(\xi_1)\bigr)
=
\varrho(\xi_1).
\]
这就得到所需序列。证毕。
\end{proof}

\noindent
上述五条结果把 H.161--H.162 的地址柱密度链与局部化数系的 prime-register 对偶链压缩为同一条有限可见性判别律：算术自然密度诱导的地址测度不是任何连续势的 Gibbs 平衡态；而在 prime-register 终对象一旦进入可数无穷层级之后，所有有限维连续观测都只能通过有限个素数轴，并必然湮灭一条余有限 \(p\)-进尾核。因而，有限半径解析证书与有限维连续表示即使联合，也仍然在解析外环与算术尾核两个方向上同时失明。
